\documentclass[12pt]{article}
\usepackage{mathtools}
\usepackage{amssymb}
\usepackage{fontspec}
\setmainfont{EB Garamond}

% Set page size and margins
\usepackage[a4paper,top=2cm,bottom=1.5cm,left=1cm,right=1cm]{geometry}
\usepackage{fancyhdr}
\pagestyle{fancy}
\setlength{\parindent}{0pt}

% Useful packages
\usepackage{minted}
\usepackage{graphicx}
\graphicspath{{./media/}}
\usepackage{subfig}
\usepackage{float}
\usepackage{siunitx}
\usepackage[hidelinks]{hyperref}
\usepackage{multirow}

% Change "Table" to "Πίνακας"
\renewcommand{\tablename}{Πίνακας}


\title{\vspace{-2cm} Θεμελιώσεις Κρυπτογραφίας \\
        Homework 5}
\author{Μαρινόπουλος Χριστόφορος 10522}
\date{Εαρινό εξάμηνο 2025}

\lhead{Homework 5}
\chead{Θεμελιώσεις Κρυπτογραφίας}
\rhead{Μαρινόπουλος Χριστόφορος}
\begin{document}
\maketitle

\section*{Ερώτημα 1}
Να εξετάσετε αν ισχύει το φαινόμενο της χιονοστιβάδας
στην συνάρτηση κατακερματισμού sha256.

\section*{Απάντηση}
Για την εξέταση του avalance effect θα εξεταστεί αν σε 100 ζευγάρια μηνυμάτων (m1,m2) που διαφέρουν σε ένα bit, η κρυπτογράφηση τους διαφέρει σε ποσοστό  >50\%. Ο κώδικας στο αρχείο \texttt{homework\_5\_10522.py} χρησιμοποιεί την βιβλιοθήκη: \texttt{pycryptodome} την οποία ο χρήστης θα πρέπει να κατεβάσει με την εντολή: \texttt{pip install pycryptodome} στο terminal εάν επιθυμεί να τρέξει τον κώδικα.
\newline
\newline
Σε αυτόν, πρώτα δημιουργείται ένα τυχαίο μήνυμα 512 bits. Μετέπειτα, δημιουργείται ένα δεύτερο το οποίο διαφέρει σε 1 bit με το πρώτο. Τέλος, κρυπτογραφούνται  και γίνεται έλεγχος εάν ο αριθμός των διαφορετικών bits των κρυπτογραφημένων μηνυμάτων είναι \( >256 \). Αυτό επαναλαμβάνεται 100 φορές και αν ο αριθμός των κρυπτογραφημένων μηνυμάτων που διαφέρουν κατά 256 bit είναι μεγαλύτερος από 50, ( περισσότερο από το 50\% του δείγματος) τότε για τον sha 256 ισχύει το Avalance effect.
\newline
\newline
Με την εκτέλεση του κώδικα παρατειρείτε πως ισχύει όντως το avalance effect.
\end{document}